% TOP_PROBLEM: {"id": 600011, "title": "Baker's Dozen — Convex tangent-segment iteration", "category_id": 6, "category": {"id": 6, "name": "geometry", "display_name": "Geometry", "description": "Euclidean and non-Euclidean geometry, geometric structures.", "slug": "geometry", "order_index": 6, "created_at": "2026-07-31T15:26:25.671Z"}, "status": "partially_solved", "problem_number": "AMR-005-0011", "set_id": 13}
% FIRST_POSTED: 2026-10-01
% ENTRY_KIND: statement_only
% UnsolvedMath ID 600011; source catalogue code AMR-005-0011
% !TeX program = lualatex
% Definition and sources only; this file is not an LLM solution attempt.
\documentclass[11pt,a4paper]{article}
\usepackage[margin=25mm]{geometry}
\usepackage{fontspec}
\setmainfont{lmroman10-regular.otf}[BoldFont=lmroman10-bold.otf,
 ItalicFont=lmroman10-italic.otf,BoldItalicFont=lmroman10-bolditalic.otf]
\usepackage{amsmath,amssymb,mathtools}
\usepackage{unicode-math}
\setmathfont{latinmodern-math.otf}
\AtBeginDocument{\let\setminus\smallsetminus}
\usepackage{xurl,hyperref}
\hypersetup{colorlinks=true,urlcolor=blue}
\setcounter{secnumdepth}{0}
\setlength{\parindent}{0pt}
\setlength{\parskip}{5pt}
\setlength{\emergencystretch}{3em}
\begin{document}
\section*{Baker's Dozen — Convex tangent-segment iteration}\label{600011}
{\small UnsolvedMath ID 600011 · AMR-005-0011 · Geometry · AMR Open Problem Lists}
\subsection{Definitions and mathematical statement}
An oriented oval is a smooth embedded closed curve $\gamma:S^1\to\mathbb R^2$
bounding a strictly convex region. For a regular parametrization define
\[
 \tau_\gamma(t)=\frac{\gamma'(t)}{\|\gamma'(t)\|},\qquad
 T\gamma(t)=\gamma(t)+\tau_\gamma(t).
\]
The image $T\gamma$ is the locus of the endpoints of oriented unit tangent
segments. Reparametrizations preserving orientation give the same geometric
curve. When iterating $T$, compute the unit tangent of the new curve each
time; one cannot keep the original speed parameter and replace all later
unit tangents by unnormalized derivatives.

Put $\gamma_0=\gamma$ and $\gamma_{j+1}=T\gamma_j$ for $j\ge0$. The rigidity
conjecture states that if every $\gamma_j$ remains an embedded convex closed
curve, then the initial curve is a Euclidean circle. Convexity concerns the
entire image and rules out folds and self-intersections.

The conclusion is up to translation, rotation, and the initial radius.
Circles do satisfy the hypothesis: the tangent offset of a circle of radius
$R$ is a circle of radius $\sqrt{R^2+1}$, with an angular shift. The problem
is whether any noncircular oval can retain convexity indefinitely.
\subsection{Short English statement}
Does indefinite preservation of convexity under unit tangent offsets
characterize circles among smooth convex closed plane curves?
\subsection{Sources}
\begin{itemize}
\item[S1] ULAM.AI and the UnsolvedMath Contributors, supplied problems.json, record 600011 (AMR-005-0011), AMR Open Problem Lists. Catalogue identity and source transcription; CC BY 4.0.
\url{https://huggingface.co/datasets/ulamai/UnsolvedMath}.
\item[S2] Tabachnikov - A Baker's Dozen of Problems (2015), Conjecture 4. Original source indexed by AMR-005-0011.
\url{https://amj.math.stonybrook.edu/html-articles/Files-2015-2024/14-01/}.
\end{itemize}
\end{document}
